\section{Generative Adversarial Networks (GANs)}
\label{sec:gan_theory}
Generative Adversarial Networks (GANs) \cite{Goodfellow_2014} are a class of generative models in which a neural network is trained to synthesize samples that resemble those drawn from a target data distribution.
GANs learn the distribution by framing training as a two-player game between a generator, which produces candidate samples, and a discriminator, which attempts to distinguish real samples from generated ones.
This adversarial formulation has proven particularly effective for high-dimensional image data, where hand-crafted parametric models are often inadequate to capture the full complexity of textures, structures, and noise statistics.
Before discussing conditional extensions and image-to-image translation frameworks, it is useful to formalize this adversarial interaction as a minimax game between generator and discriminator.

\subsection{The Minimax Game}
At the core of a GAN there are two neural networks with opposing objectives:
\begin{itemize}
    \item \textbf{Generator ($G$):} maps a latent vector $z \sim p_z(z)$ (typically a low-dimensional random variable, e.g., Gaussian or uniform) or a structured input (e.g., an image) to a synthetic sample $G(z)$ in the data space. The goal of $G$ is to produce samples that are indistinguishable from real data drawn from the true data distribution $p_{data}$.
    \item \textbf{Discriminator ($D$):} maps an input sample $x$, which may originate from either the real data distribution $p_{data}$ or the generator $G$, to a scalar $D(x) \in (0,1)$, interpreted as the probability that $x$ comes from the true data distribution $p_{data}$ rather than from the generator. The goal of $D$ is to correctly classify real versus generated samples.
\end{itemize}

The training process is formulated as a minimax two-player game with the value function $V(D, G)$ defined as:

\begin{align}
    &\text{Discriminator Objective:} & \max_D V(D, G) &= \mathbb{E}_{x \sim p_{data}(x)}[\log D(x)] + \mathbb{E}_{z \sim p_z(z)}[\log (1 - D(G(z)))] \\
    &\text{Generator Objective:} & \min_G V(D, G) &= \mathbb{E}_{z \sim p_z(z)}[\log (1 - D(G(z)))]
\end{align}

During training, $D$ is updated to maximize $V(D, G)$, i.e., to assign high scores to real samples and low scores to generated samples, while $G$ is updated to minimize $V(D, G)$, i.e., to generate samples that make $D(G(z))$ as large as possible.

In practice, different variants of this objective are often used to improve gradient properties.
A common choice is the \enquote{non-saturating} generator loss, where $G$ maximizes $\mathbb{E}_{z}[\log D(G(z))]$ instead of minimizing $\mathbb{E}_{z}[\log (1 - D(G(z)))]$, leading to stronger gradients when the discriminator is confident.
Other GAN variants (e.g., Wasserstein GANs~\cite{Arjovsky2017}, least-squares GANs~\cite{Mao2017}) modify the loss to stabilize training and improve convergence, and are widely adopted in the broader state of the art for image generation.

From a game-theoretic perspective, the ideal solution of this minimax game is a Nash equilibrium, a pair $(G^*, D^*)$ such that neither player can unilaterally improve its objective given the strategy of the other.
Under suitable assumptions and infinite model capacity, it can be shown that:
\begin{itemize}
    \item The optimal discriminator for a fixed generator is
    \begin{equation}
    D^*(x) = \frac{p_{data}(x)}{p_{data}(x) + p_g(x)},
    \end{equation}
    where $p_g$ is the distribution induced by $G$.
    \item The global optimum of the game is achieved when $p_g = p_{data}$, i.e., the generator distribution matches the true data distribution and $D^*(x) = 0.5$ for all $x$.
\end{itemize}
In this situation, the discriminator cannot distinguish real from generated samples better than random guessing, and the generator produces samples that are statistically indistinguishable from real data.
In practice, GAN training rarely reaches this ideal equilibrium due to limited model capacity, optimization difficulties, and non-convexity, but the framework provides a powerful conceptual tool for designing generative models for complex data such as radargrams.

\subsection{Image-to-Image Translation}
The original GAN formulation is \emph{unconditional}: the generator learns to map random noise $z$ to samples that follow the marginal data distribution $p_{data}(x)$, without explicit control over the content of generated outputs.
In many practical applications, however, we are interested in learning a conditional distribution $p(y \mid x)$: given an input image $x$ (e.g., in our case, a simulated radargram), we wish to generate a corresponding output image $y$ (e.g., a realistic radargram) that preserves the structural content of $x$ while matching the style or statistics of a target domain.
This task is known as \textbf{image-to-image (I2I) translation}.

Examples of I2I tasks include semantic label map $\rightarrow$ RGB image synthesis (and the inverse mapping, i.e., semantic segmentation), grayscale $\rightarrow$ colorization, low-resolution $\rightarrow$ high-resolution super-resolution, or, in this thesis, simulated RS data $\rightarrow$ real RS data.
In all these cases, the model must learn a structured mapping $G: X \to Y$ between two image domains, rather than sampling from $p_{data}(y)$ in isolation.

Conditional GANs (cGANs) extend the basic GAN framework to model such conditional distributions.
Instead of learning $p_g(x)$, a cGAN learns $p_g(y \mid x)$ by conditioning both the generator and the discriminator on the input $x$:
\begin{itemize}
    \item The generator $G$ receives $x$ (and optionally noise $z$) and outputs $G(x, z)$, an image in the target domain.
    \item The discriminator $D$ receives a pair $(x, y)$ and must decide whether $y$ is a real target image paired with $x$, or a generated sample $G(x, z)$.
\end{itemize}
The adversarial objective becomes:
\begin{equation}\label{eq:cgan_objective}
\mathcal{L}_{cGAN}(G, D) = \mathbb{E}_{x,y}[\log D(x, y)] + \mathbb{E}_{x,z}[\log (1 - D(x, G(x, z)))].
\end{equation}

Compared to the unconditional GAN objective, the discriminator now evaluates consistency between input and output, rather than realism of the output alone.
This encourages the generator to produce images that are not only visually plausible in isolation, but also semantically aligned with the conditioning input.
In practice, cGANs are often combined with additional reconstruction losses (e.g., $\ell_1$ or $\ell_2$ between $G(x)$ and $y$) to further enforce fidelity to the input structure, as done in Pix2Pix and in the supervised Sim-to-Real translation experiments of this thesis.

\subsection{Conditional GANs (cGANs)}
\label{subsec:c_gan}


In this formulation, the discriminator receives input–output pairs and learns to distinguish real pairs $(x, y)$ from synthetic pairs $(x, G(x, z))$.
The generator, in turn, is encouraged to produce outputs that are both realistic and appropriately matched to the inputs.

Pix2Pix~\cite{Isola_2017} is a canonical example of a cGAN-based I2I translation model.
It combines the adversarial loss $\mathcal{L}_{cGAN}$ with an $\ell_1$ reconstruction loss
\begin{equation}
\mathcal{L}_{\ell_1}(G) = \mathbb{E}_{x,y}\left[ \| y - G(x) \|_1 \right],
\end{equation}
yielding the overall objective
\begin{equation}
\min_G \max_D \ \mathcal{L}_{cGAN}(G, D) + \lambda \mathcal{L}_{\ell_1}(G),
\end{equation}
where $\lambda$ balances perceptual realism against faithfulness to the ground truth.
In this thesis, Pix2Pix is used as a supervised baseline to translate simulated RS radargrams into the domain of real SHARAD observations, exploiting the availability of paired examples over Elysium Planitia.

\subsection{CycleGAN}
\label{subsec:cyclegan}
For many applications, including planetary RS, paired training data $(x, y)$ are scarce or unavailable: we may have a set of simulated images $X$ and a separate set of real images $Y$, but no one-to-one correspondence between them.
CycleGAN~\cite{Zhu_2017} addresses this setting by learning two mappings $G: X \to Y$ and $F: Y \to X$ together with two discriminators $D_Y$ and $D_X$:
\begin{itemize}
    \item $D_Y$ distinguishes real images $y \in Y$ from translated images $G(x)$.
    \item $D_X$ distinguishes real images $x \in X$ from translated images $F(y)$.
\end{itemize}
Adversarial losses are defined in each domain, encouraging $G$ and $F$ to produce samples that are indistinguishable from real images in $Y$ and $X$, respectively.
To avoid arbitrary mappings that ignore the input content, CycleGAN introduces the \textbf{cycle-consistency loss}:
\begin{equation}
\mathcal{L}_{cyc}(G, F) = \mathbb{E}_{x \sim X}\left[ \| F(G(x)) - x \|_1 \right] + \mathbb{E}_{y \sim Y}\left[ \| G(F(y)) - y \|_1 \right].
\end{equation}
This loss enforces that translating an image from one domain to the other and back again should approximately recover the original, thereby preserving structural information.

The full CycleGAN objective combines adversarial and cycle-consistency terms, and optionally additional regularizers (e.g., identity loss).
In this thesis, CycleGAN is used to perform unpaired Sim-to-Real translation between simulated and real RS radargrams, enabling the generator to learn the statistics of clutter and noise in the real domain without requiring perfectly aligned simulation–observation pairs.

\subsection{Challenges in Training}
\label{subsec:gan_challenges}
GAN training is notoriously unstable due to the non-convex, high-dimensional optimization landscape and the competing objectives of the two networks.
Several failure modes are commonly observed, particularly when training on complex data such as radargrams.
Below we describe the main challenges and established mitigation strategies.

\begin{itemize}
    \item \textbf{Mode Collapse:} The generator collapses to producing only a limited subset of the data distribution (e.g., a single \enquote{average} radargram texture), fooling the discriminator but failing to capture the full diversity of clutter patterns, stratigraphy, and noise statistics in RS data. This is particularly problematic for high-dimensional images where the generator can easily overfit to the most salient samples.

    Possible mitigation techniques are:
    \begin{itemize}
        \item \textbf{Spectral normalization} or \textbf{gradient penalty} on the discriminator to enforce Lipschitz continuity and prevent it from becoming too powerful.
        \item \textbf{Mini-batch discrimination} using statistics (e.g., variance) rather than per-sample classification.
        \item \textbf{Unrolled GANs} or \textbf{progressive growing}, where the discriminator is trained for multiple steps per generator update.
        \item \textbf{Auxiliary losses} such as feature matching, where the generator is encouraged to match intermediate features of the discriminator rather than its final output.
    \end{itemize}

    \item \textbf{Vanishing Gradients:} When the discriminator becomes too confident ($D(G(z)) \approx 0$), the generator receives vanishingly small gradients from the standard loss $\log(1 - D(G(z)))$, halting learning.

    Possible mitigation techniques are:
    \begin{itemize}
        \item \textbf{Non-saturating loss} for the generator: maximize $\log D(G(z))$ instead of minimizing $\log(1 - D(G(z)))$.
        \item \textbf{Label smoothing}, where real samples are assigned soft labels (e.g., $D(x) = 0.9$ instead of $1.0$) to prevent overconfidence.
        \item \textbf{Noise injection} in the discriminator input or \textbf{dropout} to maintain gradient flow.
    \end{itemize}

    \item \textbf{Training Imbalance:} The discriminator can overpower the generator early in training, or vice versa, leading to oscillations or collapse.

    Possible mitigation techniques are:
    \begin{itemize}
        \item \textbf{Adaptive training schedules}: update the weaker network more frequently or use heuristics to pause the stronger one.
        \item \textbf{Wasserstein GANs (WGAN)} \cite{Arjovsky2017}, which replace the Jensen–Shannon divergence with the Wasserstein distance and enforce $1$-Lipschitz continuity via gradient penalty.
        \item \textbf{Two-timescale update rule (TTUR)} \cite{Heusel2017}, where the generator and discriminator have different learning rates and optimizers.
    \end{itemize}

    \item \textbf{Representation Collapse:} The generator learns a representation that is statistically similar to the data but lacks fine semantic details (e.g., blurry radargrams that match average speckle statistics but fail to preserve reflector geometry).

    Possible mitigation techniques are:
    \begin{itemize}
        \item \textbf{Perceptual losses} such as LPIPS or feature matching.
        \item \textbf{Multi-scale discriminators} that operate at different resolutions.
        \item \textbf{Progressive growing of GANs (ProGAN)} \cite{Karras2018}, where the model is trained progressively from low to high resolution.
    \end{itemize}
\end{itemize}

In this thesis, we adopt a combination of Instance Normalization (in the discriminator), Least Squares loss (LSGAN) (for the generator), Label Smoothing, and mixed precision training to stabilize GAN training on RS data.
While more advanced techniques (e.g., WGAN-GP or Spectral Normalization) could be explored in future work, this baseline setup yields robust convergence for both Pix2Pix and CycleGAN on the SHARAD dataset.

A detailed discussion of the evaluation metrics used to assess generative quality in the experimental section of this thesis, including pixel-wise (PSNR, SSIM), perceptual (LPIPS), and domain-specific measures, is provided in Section~\ref{subsec:evaluation_metrics}.

